\section{Proofs}

\subsection{Consumer access and channel retention}

\begin{lemma}[Consumer access]
\label{lem:consumer}
Under equations \eqref{eq:consumer_ces} and \eqref{eq:platform_price},
\[
s'(z)=(\eta-1)s(z)[1-s(z)]>0
\quad\text{and}\quad
\frac{d\ln P_M(z)}{dz}=-s(z)<0.
\]
If \(\ell_A>\ell_P\), then \(\omega_C'(z)<0\) and
\(\Lambda(z)>0\).
\end{lemma}

\begin{proof}
Write:
\[
x(z)=a_Pp_P(z)^{1-\eta},
\qquad
a=(1-a_P)p_A^{1-\eta}.
\]
Then \(s=x/(a+x)\).  Because \(d\ln p_P/dz=-1\),
\[
\frac{x'(z)}{x(z)}=-(1-\eta)=\eta-1.
\]
Therefore:
\[
s'
=\frac{x'a}{(a+x)^2}
=(\eta-1)\frac{x}{a+x}\frac{a}{a+x}
=(\eta-1)s(1-s)>0.
\]
For the price index:
\[
\ln P_M=\frac{1}{1-\eta}\ln(a+x).
\]
Hence:
\[
\frac{d\ln P_M}{dz}
=\frac{1}{1-\eta}\frac{x'}{a+x}
=\frac{\eta-1}{1-\eta}s
=-s.
\]
Finally:
\[
\omega_C'(z)=(\ell_P-\ell_A)s'(z).
\]
It is negative when \(\ell_A>\ell_P\).  Since \(D>0\), the definition in
equation~\eqref{eq:capture_drag} then implies \(\Lambda>0\).
\end{proof}

At the limits \(s\rightarrow0\) and \(s\rightarrow1\), \(s'\rightarrow0\).
The platform continues to reduce \(P_M\), but the expenditure-reallocation
margin vanishes at either channel boundary.  This is why capture drag is
largest at an intermediate platform share rather than at complete platform
penetration.

\subsection{Endowment-conditioned production costs}

The background technology in equation~\eqref{eq:background_technology}
implies:
\[
w_i=(1-\theta_0)\frac{Y_{i0}}{L_i},
\qquad
r_i=\theta_0\frac{Y_{i0}}{K_i}.
\]
With \(Y_{i0}=A_0K_i^{\theta_0}L_i^{1-\theta_0}\), division by the relevant
factor endowment yields equation~\eqref{eq:factor_prices}.  Substituting those
factor prices into the candidate industry's Cobb--Douglas unit expenditure
function gives:
\begin{align*}
c^P_{ij}
&=
\frac{1}{A_j}
\left(\frac{\theta_0A_0e_i^{\theta_0-1}}{\theta_j}\right)^{\theta_j}
\left(
\frac{(1-\theta_0)A_0e_i^{\theta_0}}{1-\theta_j}
\right)^{1-\theta_j}\\
&=
\frac{A_0}{A_j}
\left(\frac{\theta_0}{\theta_j}\right)^{\theta_j}
\left(\frac{1-\theta_0}{1-\theta_j}\right)^{1-\theta_j}
e_i^{\theta_0-\theta_j}\\
&\equiv \xi_j e_i^{\theta_0-\theta_j}.
\end{align*}
Thus:
\[
\frac{\partial\ln c^P_{ij}}{\partial\ln e_i}
=\theta_0-\theta_j.
\]
For \(\theta_j>\theta_0\), capital abundance lowers the candidate industry's
production cost relative to the background opportunity cost.  LCA still
requires the external-reference double ratio in equation~\eqref{eq:lca};
the derivative is only the local structural component of that comparison.

\subsection{Proof of Proposition~\ref{prop:entry}}

\begin{proof}
For a given organization \(r\), equation~\eqref{eq:revenue_profit} becomes:
\[
\pi_r(z,G)
=
k\Mcal_0e^{\varepsilon z+\Omega(z)G}(c+d_r)^{1-\sigma}
-F_r.
\]
It is strictly increasing in \(G\) whenever \(\Omega(z)>0\).  The unique raw
zero-profit threshold is equation~\eqref{eq:mode_entry_threshold}.  Direct
differentiation gives:
\[
\frac{\partial G_r^0}{\partial c}
=
\frac{\sigma-1}{\Omega(z)(c+d_r)}>0.
\]
The minimum of two strictly increasing functions is strictly increasing:
for \(c_2>c_1\), both \(G_P^0(c_2)>G_P^0(c_1)\) and
\(G_L^0(c_2)>G_L^0(c_1)\), so their minimum at \(c_2\) exceeds their minimum
at \(c_1\).  The truncation at zero preserves weak monotonicity.  At an
interior threshold, the relationship is strict.

Holding the external-region and background-industry ratios fixed,
\(\partial c/\partial\mathcal C^P_{ij}>0\).  The chain rule therefore gives:
\[
\frac{\partial G_E}{\partial\mathcal C^P_{ij}}
=
\frac{\partial G_E}{\partial c}
\frac{\partial c}{\partial\mathcal C^P_{ij}}>0
\]
at an interior threshold.

For the ACA interpretation, the focal region's delivered marginal cost is
\(\tau_O(z,G)(c+d_r)\).  Since
\[
\frac{\partial\ln\tau_O(z,G)}{\partial G}
=-(\gamma+\chi z)<0,
\]
infrastructure lowers the focal region's total-cost ratio relative to a fixed
reference region.  Stronger LCA starts from a lower production-cost ratio, so
it requires weakly less \(G\) both to achieve a total-cost advantage and to
cover the fixed cost.  The first condition is ACA; the second is viability.
Entry and sales are realization.
\end{proof}

\section{Complete Entry and Organization Classification}

Let:
\[
b_P=k(c+d_P)^{1-\sigma},
\qquad
b_L=k(c+d_L)^{1-\sigma},
\]
so \(b_L>b_P>0\).  It is convenient to express the state problem in terms of
effective market scale:
\[
X(z,G)=\Mcal_0e^{\varepsilon z+\Omega(z)G}.
\]
Profits are:
\[
\pi_P=Xb_P-f,
\qquad
\pi_L=Xb_L-f-F.
\]
The corresponding market-scale thresholds are:
\[
X_P=\frac{f}{b_P},
\qquad
X_{L0}=\frac{f+F}{b_L},
\qquad
X_X=\frac{F}{b_L-b_P}.
\]

\begin{lemma}[Three-region condition]
\label{lem:three_region}
Define:
\[
\bar F=f\frac{b_L-b_P}{b_P}.
\]
If \(F>\bar F\), then:
\[
X_P<X_{L0}<X_X,
\]
equivalently \(G_P^0<G_L^0<G_X\), and the equilibrium states are:
\[
\begin{cases}
0,&X<X_P,\\
P,&X_P\leq X<X_X,\\
L,&X\geq X_X.
\end{cases}
\]
If \(F=\bar F\), all three thresholds coincide.  If \(F<\bar F\), the
platform-dependent interval is absent and the producer enters through \(L\).
\end{lemma}

\begin{proof}
First:
\begin{align*}
X_P<X_{L0}
&\Longleftrightarrow
\frac{f}{b_P}<\frac{f+F}{b_L}\\
&\Longleftrightarrow
f(b_L-b_P)<Fb_P\\
&\Longleftrightarrow F>\bar F.
\end{align*}
Second:
\begin{align*}
X_{L0}<X_X
&\Longleftrightarrow
\frac{f+F}{b_L}<\frac{F}{b_L-b_P}\\
&\Longleftrightarrow
f(b_L-b_P)<Fb_P\\
&\Longleftrightarrow F>\bar F.
\end{align*}
Thus the same condition orders all three thresholds.  Below \(X_P\), both
profits are negative.  At \(X_P\), \(\pi_P=0\), while:
\[
\pi_L-\pi_P
=X_P(b_L-b_P)-F
=\bar F-F<0.
\]
Because:
\[
\pi_L-\pi_P=X(b_L-b_P)-F
\]
is strictly increasing in \(X\), it crosses zero exactly once at \(X_X\).
This proves the three states.  Equality collapses all thresholds to one
point.  If \(F<\bar F\), the local mode is already more profitable at the
first viable entry scale, so \(P\) is never selected.
\end{proof}

\section{Proof of Proposition~\ref{prop:embedding}}

\begin{proof}
Define:
\[
\Delta(c)
=(c+d_L)^{1-\sigma}-(c+d_P)^{1-\sigma}>0.
\]
The raw switch threshold is:
\[
G_X
=
\frac{\ln F-\ln[k\Mcal_0e^{\varepsilon z}]-\ln\Delta(c)}
{\Omega(z)}.
\]
Its cost derivative is:
\[
\frac{\partial G_X}{\partial c}
=-\frac{\Delta'(c)}{\Omega(z)\Delta(c)}.
\]
Because:
\[
\Delta'(c)
=(1-\sigma)
\left[(c+d_L)^{-\sigma}-(c+d_P)^{-\sigma}\right]<0,
\]
we have \(\partial G_X/\partial c>0\).  Truncation at zero preserves weak
monotonicity.  The LCA result follows from
\(\partial c/\partial\mathcal C^P_{ij}>0\).

For \(G>G_X\):
\[
\pi_L-\pi_P
=
k\Mcal_0e^{\varepsilon z+\Omega G}\Delta(c)-F>0.
\]
Hence \(L\) is selected whenever it also weakly dominates no entry; the
classification in Lemma~\ref{lem:three_region} covers the possible entry
orders.

Since \(d_L<d_P\):
\[
(c+d_L)^{1-\sigma}>(c+d_P)^{1-\sigma},
\]
and therefore \(\Rcal_L>\Rcal_P\) at a common \((z,G)\).  The payment-flow
definitions imply:
\[
\rho_L=1>
1-\frac{d_P}{\mu(c+d_P)}
=\rho_P.
\]
Both inequalities are strict, so:
\[
B_j^L=\rho_L\Rcal_L
>
\rho_P\Rcal_P=B_j^P.
\]
The external payment changes from \(d_Pq_P>0\) to zero in the polar baseline,
while the unit organizational service \(d_Lq_L\) and the shared fixed
organizational service \(F\) are locally supplied.  The organizational switch
therefore changes local service income, external payments, and total
first-round local income.
\end{proof}

\section{Proof of Proposition~\ref{prop:income_thresholds}}

\subsection{Income derivatives within a producer state}

Fix an active producer organization \(r\).  The local producer-income term can
be written:
\[
B_j^r(z,G)
=
A_r
\exp\{\varepsilon z+\Omega(z)G\},
\]
where:
\[
A_r
=
\rho_r\Mcal_0\mu^{1-\sigma}(c+d_r)^{1-\sigma}>0.
\]
Therefore:
\[
\frac{\partial\ln B_j^r}{\partial G}
=\Omega(z)>0,
\]
and:
\[
\frac{\partial\theta_r}{\partial G}
=
\Omega(z)\theta_r(1-\theta_r)>0.
\]
Define:
\[
a(G)=\varepsilon+(\sigma-1)\chi G\geq0.
\]
The nominal integration effect within organization \(r\) is:
\[
h_Y(G)=\theta_r(G)a(G)-\Lambda(z).
\]
Its derivative is:
\begin{equation}
h_Y'(G)
=
\Omega(z)\theta_r(1-\theta_r)a(G)
+\theta_r(\sigma-1)\chi>0.
\label{eq:nominal_effect_derivative}
\end{equation}
The price term \(\alpha_Ms(z)\) is independent of \(G\), so:
\[
h_R'(G)=h_Y'(G)>0
\]
within an active state.  In the no-entry state, \(\theta=0\); both effects
are constant in \(G\) until entry.

\subsection{Jumps at state changes}

Immediately before entry, \(\theta=0\), so \(h_Y=-\Lambda\).  Immediately
after entry, the viable mode generates positive sales and local producer
income, so \(\theta>0\).  Since \(a(G)\geq0\), the integration effect jumps
weakly upward and strictly upward if \(a(G)>0\).

At a \(P\)-to-\(L\) switch, Proposition~\ref{prop:embedding} gives
\(B_j^L>B_j^P\).  Hence \(\theta_L>\theta_P\), while \(a(G)\) and
\(\Lambda(z)\) are unchanged at the switch.  Both nominal and real effects
jump strictly upward.

\subsection{Existence, uniqueness, and ordering}

\begin{proof}[Completion of the proof]
The preceding derivatives and jump results imply that the integration effects
are globally nondecreasing in \(G\), strictly increasing in every active
organization, and have no downward state jumps.  Under the baseline
\(\ell_A>\ell_P\) and an interior platform share, the no-entry effect is
strictly negative because \(\Lambda>0\).

As \(G\rightarrow\infty\), the local mode is selected,
\(\theta_L\rightarrow1\), and:
\[
a(G)=\varepsilon+(\sigma-1)\chi G\rightarrow\infty
\]
when \(\chi>0\).  Hence \(h_Y(G)\rightarrow\infty\), and the real effect is
larger still.  More generally, it is sufficient to assume a finite upper
endpoint at which the relevant effect is positive.  Monotonicity and endpoint
crossing therefore imply a unique minimum threshold, allowing the crossing to
occur either continuously or at an upward jump.

Pointwise:
\[
h_R(G)=h_Y(G)+\alpha_Ms(z)>h_Y(G).
\]
Thus every \(G\) that makes the nominal effect nonnegative also makes the real
effect positive, establishing \(G_R\leq G_Y\).

Suppose both roots lie in the same active continuous organization and are
interior.  At \(G_Y\), \(h_Y(G_Y)=0\), so:
\[
h_R(G_Y)=\alpha_Ms(z)>0.
\]
Since \(h_R\) is continuous and strictly increasing in that organization, its
zero must occur strictly below \(G_Y\).  Hence \(G_R<G_Y\).  If both effects
jump from negative to nonnegative at the same discrete state change, their
minimum thresholds are identical even though \(h_R>h_Y\) on either side.

It remains to establish the comparative static with respect to LCA.  For the
local mode:
\[
B_j^L
=
\Mcal\mu^{1-\sigma}(c+d_L)^{1-\sigma},
\]
which is strictly decreasing in \(c\).  For the platform mode:
\begin{align*}
B_j^P
&=
\Mcal\mu^{1-\sigma}
\left[
(c+d_P)^{1-\sigma}
-\frac{d_P}{\mu}(c+d_P)^{-\sigma}
\right].
\end{align*}
Differentiating the bracket gives:
\begin{align*}
&(1-\sigma)(c+d_P)^{-\sigma}
+\frac{\sigma d_P}{\mu}(c+d_P)^{-\sigma-1}\\
&\qquad
=-(\sigma-1)c(c+d_P)^{-\sigma-1}<0,
\end{align*}
where \(\mu=\sigma/(\sigma-1)\).  Thus lower \(c\) raises local producer
income in either mode.  It also weakly advances entry and the switch to the
higher-income local mode by Propositions~\ref{prop:entry} and
\ref{prop:embedding}.  At every fixed \(G\), stronger LCA therefore weakly
raises \(\theta_{r^*}\) and the two integration effects.  Their minimum
thresholds weakly fall.  The relationship is strict when the relevant root
remains interior to the same active organizational state.
\end{proof}

\section{Proof of Proposition~\ref{prop:state}}

\begin{proof}
Define the infrastructure provider's first-order condition:
\[
\Phi(G,\vartheta)
=
\mathcal U_I'(G)-\kappa_I(\vartheta)C'(G)=0.
\]
Its derivative with respect to \(G\) is:
\[
\Phi_G
=
\mathcal U_I''(G)-\kappa_I(\vartheta)C''(G)<0.
\]
The objective is therefore strictly concave and has at most one interior
solution.  Under the maintained endpoint conditions, that solution exists and
is unique.  Since:
\[
\Phi_\vartheta
=-\kappa_I'(\vartheta)C'(G)>0,
\]
implicit differentiation gives:
\[
G'(\vartheta)
=-\frac{\Phi_\vartheta}{\Phi_G}>0.
\]

For any fixed structural threshold \(G_H\), the monotonicity of
\(G(\vartheta)\) implies that the set
\(\{\vartheta:G(\vartheta)\geq G_H\}\) is an interval.  When nonempty, it has
a unique minimum \(\vartheta_H\), except for an economically irrelevant flat
boundary if the infrastructure solution itself is constrained.  Increasing
\(\vartheta\) reduces \(G_H-G(\vartheta)\).  Stronger LCA lowers
\(G_E,G_L,G_R,G_Y\) by Propositions~\ref{prop:entry}--\ref{prop:income_thresholds};
applying the increasing inverse of \(G(\vartheta)\) weakly lowers the
associated participation thresholds.

For the protection comparison, a consumer wedge \(\tau_C>1\) changes the
effective platform price to \(\tau_C\bar p_Pe^{-z}\).  The CES price index is
increasing in an input price, so protection raises \(P_M\) and lowers the
platform expenditure share.  On the producer side:
\[
\frac{\partial G_P^0}{\partial d_P}
=
\frac{\sigma-1}{\Omega(z)(c+d_P)}>0.
\]
Thus protection makes the low-fixed-cost external entry mode harder for a
marginal producer.

Let:
\[
\Delta=(c+d_L)^{1-\sigma}-(c+d_P)^{1-\sigma}.
\]
Because:
\[
\frac{\partial\Delta}{\partial d_P}
=
(\sigma-1)(c+d_P)^{-\sigma}>0,
\]
the switch threshold satisfies:
\[
\frac{\partial G_X}{\partial d_P}
=
-\frac{1}{\Omega(z)\Delta}
\frac{\partial\Delta}{\partial d_P}<0.
\]
Protection can therefore induce a switch toward \(L\), but it does so by
worsening \(P\), while simultaneously harming consumer access and making
platform-based entry more difficult.  Facilitation raises \(G\), which
increases market access under both modes and does not enter the consumer price
equation.  The policies are not equivalent.
\end{proof}

\section{Boundary Cases and Counterexamples}

\subsection{No consumer-side capture gap}

If \(\ell_A=\ell_P\), then:
\[
\omega_C(z)=\ell_A,
\qquad
\Lambda(z)=0.
\]
The nominal effect becomes:
\[
\frac{d\ln Y}{dz}
=
\theta_r[\varepsilon+(\sigma-1)\chi G]\geq0
\]
in every active producer state.  A negative nominal effect can then arise
only in an extension with another first-round loss.  This case confirms that
the baseline negative region is generated by an explicit payment-incidence
gap rather than platform terminology.

If \(\ell_P>\ell_A\), then \(\Lambda<0\).  Platform substitution raises
consumer-side retention as well as reducing the price index.  Both nominal
and real effects shift upward; the low-infrastructure loss region may vanish.

\subsection{No producer-side spatial payment gap}

Setting \(\lambda=1\) in equation~\eqref{eq:partial_localization} gives
\(\rho_P=1=\rho_L\).  The organization switch can still raise sales because
\(d_L<d_P\), but it no longer changes the spatial retention rate.  This case
separates the cost advantage of local organization from its payment-location
advantage.

\subsection{No infrastructure--platform complementarity}

If \(\chi=0\), producer market access is:
\[
\Mcal(z,G)
=
\Mcal_0e^{\varepsilon z+(\sigma-1)\gamma G}.
\]
Infrastructure can still raise the level of candidate-industry income and
therefore \(\theta_r\), but it no longer increases the marginal producer
benefit of platform integration directly.  Equation~\eqref{eq:nominal_effect}
loses the term \((\sigma-1)\chi G\).  The threshold mechanism can survive,
but the strong complementarity interpretation cannot.

\subsection{Collapse of the platform-dependent region}

Lemma~\ref{lem:three_region} gives the exact counterexample to a universal
three-stage development path.  When \(F\leq\bar F\), the local mode is
sufficiently inexpensive in fixed-cost terms that the candidate industry
enters directly through \(L\).  Platform dependence is therefore neither
necessary nor universal.

\subsection{Weak rather than strict income-threshold ordering}

Suppose an entry or organizational switch occurs at \(\widehat G\), with:
\[
h_R(\widehat G^-)<0,\qquad h_Y(\widehat G^-)<0,
\]
and:
\[
h_R(\widehat G^+)\geq0,\qquad h_Y(\widehat G^+)\geq0.
\]
Then:
\[
G_R=G_Y=\widehat G.
\]
The pointwise inequality \(h_R>h_Y\) is not enough to create two different
minimum thresholds when a common discrete jump crosses both zeroes.  This is
why Proposition~\ref{prop:income_thresholds} states strict ordering only under
a continuous interior-root condition.

\subsection{Instability of the local-income closure}

If:
\[
D(z)=1-\beta-\alpha_M\omega_C(z)\leq0,
\]
the household-spending recursion does not converge.  This is not an
alternative equilibrium of the paper.  It is an inadmissible parameter
region.  The baseline values satisfy \(D(z)>0\) for the full numerical
support.

\section{Mapping to the Companion EL Income Block}

The companion access--capture model can be written:
\[
Y^{EL}
=
\frac{B^{EL}+\omega^{EL}G^{EL}}
{1-\beta^{EL}-\alpha^{EL}\omega^{EL}}.
\]
The present paper does not insert this expression into the producer problem.
It uses the following one-to-one mapping only after first-round payments have
been determined:

\begin{table}[htbp]
\centering
\caption{Mapping to the companion income-propagation block}
\begin{tabular}{p{0.28\textwidth}p{0.29\textwidth}p{0.35\textwidth}}
\toprule
Companion object & Present object & Accounting rule\\
\midrule
\(B^{EL}+\omega^{EL}G^{EL}\)
& \(B_0+B_j^{r^*}\)
& autonomous first-round local income only\\
\(\alpha^{EL}\)
& \(\alpha_M\)
& household market-good expenditure share\\
\(\beta^{EL}\)
& \(\beta\)
& locally supplied nontraded-good expenditure share\\
\(\omega^{EL}\)
& \(\omega_C(z)\)
& local income content of induced market-good spending\\
consumer price index
& \(P_M(z)\)
& identical CES market-good price concept\\
nominal income
& \(Y\)
& equation~\eqref{eq:nominal_income}\\
consumption-equivalent income
& \(R\)
& equation~\eqref{eq:real_income}\\
\bottomrule
\end{tabular}
\end{table}

The candidate-industry sales payment enters \(B_j^{r^*}\) once.  It does not
also enter \(\omega_C\), which applies only to subsequent household
market-good expenditure.  The external platform organization payment enters
\(\rho_P\) once.  Consumer-channel external payments enter \(\ell_P\) and
\(\omega_C\), not \(\rho_P\).  This separation prevents production,
organizational services, and induced household spending from being counted
twice.

\section{Numerical Reproducibility and Mechanism Closures}

The numerical implementation consists of:
\begin{itemize}
 \item \texttt{code/numerical/model.py}, which implements the equations and
 equilibrium choices;
 \item \texttt{code/numerical/generate\_figures.py}, which generates all
 figures and source-data CSV files;
 \item the validation script, which checks analytical derivatives, finite
 differences, threshold uniqueness, the three-region condition, boundary
 cases, and state monotonicity.
\end{itemize}

For \(z=0.5,c=0.8\), the code obtains:
\[
G_P^0=-0.070176,\quad
G_L^0=0.126129,\quad
G_X=0.256065,
\]
\[
G_R=0.046348,\quad
G_Y=0.550487.
\]
Finite-difference derivatives have the signs established in the proofs.
The program also verifies that lowering \(F\) below \(\bar F\) removes the
platform-dependent interval and that equal consumer-channel local contents
set \(\Lambda\) exactly to zero.

The file
\texttt{figures/source\_data/appendix\_mechanism\_closures.csv}
reports the full \(G\)-grid for:
\begin{enumerate}
 \item the baseline;
 \item \(\ell_A=\ell_P\);
 \item complete platform-service localization;
 \item \(\chi=0\);
 \item partial platform-service localization.
\end{enumerate}
All reported exercises use normalized, illustrative values.  No parameter is
estimated, and the exercises do not constitute a calibration.
